% An echo-type sequence applied to the root pulse partially mitigates decoherence effects but remains limited by pronounced amplitude-dependent broadening. Conversely, the Lorentzian echo variant demonstrates an amplitude-independent linewidth spanning over three orders of magnitude in drive strength, evidencing remarkable robustness against amplitude-induced dephasing. This insensitivity arises from the adiabatic evolution and the symmetric frequency-splitting intrinsic to the pulse design. The primary limitation lies in the floating offset and amplitude-dependent baseline, resulting in reduced contrast compared to conventional square and Lorentzian drives.



To identify the useful operating regime, we examine the linewidth and contrast
of the echo response as functions of pulse duration and drive amplitude.  We
choose pulse durations shorter than $T_1$, but long enough that Fourier
broadening does not dominate.  The duration-dependent linewidth, contrast,
and contrast-weighted resolution are reported in the Supplemental
Material~\cite{SupplementalMaterial}.

A second design constraint is the cutoff of the formally infinite-support
pulse envelope.  Truncation introduces a residual endpoint drive
$\Omega_c=c\Omega_0$ and a corresponding cutoff-broadening
contribution~\cite{boradjiev_qubit_2013,mihov_defying_2024,mihov_qubit_2024}.
The relevant weak-saturation condition is
$s_c=\Omega_c^2T_1T_2\ll1$, or
$\Omega_c\ll1/\sqrt{T_1T_2}$; the inequality therefore points toward
\emph{smaller}, not larger, endpoint amplitudes.  The relative cutoff $c$
cannot be specified independently of $\Omega_0$, $T_1$, and $T_2$, giving a
three-dimensional design space of pulse duration, drive amplitude, and cutoff
level.  The corresponding cutoff--amplitude measurements and their
resolution--signal analysis are reported in the Supplemental
Material~\cite{SupplementalMaterial}.

Beyond linewidth and contrast, amplitude robustness requires that increasing
the drive not displace the inferred resonance center.  At large amplitudes,
multilevel AC-Stark shifts and finite-pulse interference become important; we
therefore examine the center stability in an additional q6 measurement
over a calibrated nominal Rabi-frequency range of $10$--$60~\mathrm{MHz}$.

\begin{figure}[t]
    \centering
    \includegraphics[width=\columnwidth]{%
        ../figures/paper/04_main_ac_stark_correction_maps.pdf}
    \caption{Measured q6 echo-root-Lorentzian spectroscopy for (a) $\beta=0$
    and (b) $\beta=-0.22$, with $L=50~\mu\mathrm{s}$ and cutoff $c=0.00075$.
    White curves trace the fitted resonance centers. Over the
    $10$--$60~\mathrm{MHz}$ drive range, the $\beta=-0.22$ waveform shows
    smaller amplitude-dependent fluctuations of the resonance center,
    while preserving the narrow depletion feature. This waveform combines
    a DRAG quadrature with a smoothed midpoint inversion.}
    \label{fig:ac-stark-correction-maps}
\end{figure}

\begin{figure}[t]
    \centering
    \includegraphics[width=\columnwidth]{%
        ../figures/paper/04_main_ac_stark_shifts_square.pdf}
    \caption{Resonance-center stability and linewidth of the measured
    echo-root-Lorentzian response. (a) Measured center variations
    compared with the simulated constant-drive AC-Stark shift. The measured
    curves use the same fitted centers as Fig.~\ref{fig:ac-stark-correction-maps}
    for $\beta=0$ and $-0.22$. The inset resolves their amplitude-dependent
    fluctuations over $10$--$60~\mathrm{MHz}$: the $\beta=-0.22$ waveform
    exhibits smaller center variations. The constant-drive curve retains
    the three-level simulation with
    $\alpha/(2\pi)=\MeasuredAnharmonicity~\mathrm{MHz}$.
    (b) Gaussian-fit FWHM of the same measured spectra and the red
    constant-drive power-broadening reference, on a logarithmic scale.
    The right axis expresses linewidths in units of the paper's q1
    reference $\Gamma_{T_2}=\EffectiveCoherenceFwhm~\mathrm{kHz}$ (dotted line).
    Both measured waveforms broaden with drive, while the $\beta=0$
    response exhibits larger linewidth excursions at high amplitudes.}
    \label{fig:main-ac-stark-shifts}
\end{figure}

For comparison, the weak-drive three-level AC-Stark shift of a constant drive
is $\delta f_{01}\simeq-\left(\Omega_0/2\pi\right)^2/
[2(\alpha/2\pi)]$~\cite{PhysRevA.76.042319,schneider_multilevel_2018}.
With $\alpha/(2\pi)=\MeasuredAnharmonicity~\mathrm{MHz}$, it reaches
$5.8~\mathrm{MHz}$ at $\Omega_0/(2\pi)=50~\mathrm{MHz}$.  For the
$20~\mu\mathrm{s}$ root-Lorentzian envelope with $c=0.001$, the phase-average
shift is smaller by $\langle f^2\rangle_t=0.001570$ and reaches only
$9.1~\mathrm{kHz}$; the echo sign change does not cancel this quadratic term.
Fig.~\ref{fig:ac-stark-correction-maps} instead tests a derivative-quadrature
waveform on q6 in the September 5--6 calibration campaign. The acquisition
uses $I+iQ$ with $Q=-\beta\dot I/(2\pi\alpha_{\rm cal})$, where
$\alpha_{\rm cal}=237.95~\mathrm{MHz}$ is the signed value supplied to the
waveform generator and time is in microseconds. Both the derivative
quadrature and the midpoint smoothing change between these sequential runs;
no accumulated-phase correction is applied. The depletion feature persists
across the displayed $10$--$60~\mathrm{MHz}$ range in both maps. For display,
we subtract each run's arithmetic mean accepted fitted center over this
range, $20.85$ and $36.68~\mathrm{kHz}$ relative to the acquisition reference
for $\beta=0$ and $-0.22$, respectively. The RMS residual center motion is
$6.78$ and $2.64~\mathrm{kHz}$ about these respective means. This lower
scatter characterizes the combined smoothed-midpoint and DRAG waveform;
the comparison does not isolate the derivative-quadrature contribution.
The absence of resolved three-state readout precludes a measured
leakage-suppression claim.

Fig.~\ref{fig:main-ac-stark-shifts}(a) compares these measured center
variations with the original three-level constant-drive simulation. The
experimental curves retain the centering used in
Fig.~\ref{fig:ac-stark-correction-maps}; the simulated curve is referenced
to the bare transition and retains the q1 anharmonicity
$\alpha/(2\pi)=\MeasuredAnharmonicity~\mathrm{MHz}$ as an illustrative scale
comparison. Its megahertz-scale displacement contrasts with the kilohertz-scale
residual motion of the measured echo resonances. The inset magnifies the
$\pm16~\mathrm{kHz}$ region around the mean centers.
Panel (b) compares FWHM values extracted
from the same Gaussian dip estimator with the steady-state Bloch reference
$\Gamma=\Gamma_{T_2}\sqrt{1+\Omega^2T_1T_2}$, using the paper's q1 values
$T_1=\MeasuredTOne~\mu\mathrm{s}$ and $T_2=\EffectiveTTwo~\mu\mathrm{s}$.
The right axis uses the same q1 coherence reference, rather than a separate
q6 coherence calibration. Both measured linewidths increase with drive;
the smaller center fluctuations for $\beta=-0.22$ therefore do not imply
a uniformly narrower resonance. Numerical accumulated-phase compensation
is discussed separately in the Supplemental Material~\cite{SupplementalMaterial}.
